\section{Appendix}


We wanted to define braided operads as functors with certain properties over a suitable category. 
A first intuition led us to the assumption that colored braided operads might 
have a functor into the crossed simplicial braid group as defined in \cite{Dyckerhoff} such that
it fibered over inert morphisms and satisfied a weak version of a Segal condition.
However, we discovered that we were too quick to make this assumption. 
We will elaborate in the following why this did not work out as we hoped originally.

\begin{defin}\label{crossed_braid} The {\em crossed simplicial braid group} is a category $\Delta B$ equipped with an embedding
	$\iota:\Delta \to \Delta B$ such that:
	\begin{enumerate}
		\item the functor $\iota$ is bijective on objects,
		\item any morphism $u: \iota[m] \to \iota[n]$ in $\Delta B$ can be uniquely expressed as 
			a composition $\iota(\phi) \circ b$ where $\phi: [m] \to [n]$ is a morphism in
			$\Delta$ and $b \in B_m$ is acting on $\iota[m]$ in $\Delta B$.
	\end{enumerate}
    Note that we sometimes refer to $\iota[m]$ by writing $[m]$ instead.
\end{defin}

\begin{rem}
    To understand composition in the category $\Delta B$ let us consider two morphisms
    \[
    f \colon [m] \to [n] \quad \text{and} \quad g \colon [n] \to [l]
    \]
    with the factorization $f = \phi_f \circ b_f$ and $g = \phi_g \circ b_g$. 
    We define for every $j \in [n]$ an $f$-pencil $p_j$ as a collection $[j, i_1, \dots i_s]$ 
    with $\phi_f^{-1}(j) = \{i_1 < \dots < i_s\}$. We call $s$ the width of the pencil. Note 
    that we say it is an empty pencil $p_j$ if $\phi_f^{-1}(j) = \emptyset$.
    
    Then we obtain a braid on $m$ strands in the following way:
    First, we consider $p_1, \dots, p_n$ as $n$ strand such that we can braid by $b_g$. Then, 
    we forget the empty pencil strands; in other words we restrict the action of $b_g$ to the 
    subset $J \subset [n]$ such that $p_j$ is not empty for all $j \in J$. However, for notation reason and without loss of generality, we assume that there are no empty pencils.    
    We obtain a braiding on $m$ strands if we consider the action of $b_g$ on the pencils as a 
    block braiding on $m$ strand. Each pencil corresponds to a block having the length of the 
    pencils width. The blocks are ordered by $[n]$.

    Geometrically this can be thought of as braiding the pencils. Then, as the order within the pencil remains unchanged, we can cut a pencil $[j, j_1, \dots, j_s]$ into its $s$ underlying strands. The resulting braid corresponds to $b_g(s_1, \dots, s_n) \circ b_f$ where $s_j$ is the width of the pencil $p_j$ for $j \in [n]$.
    Thus, we define
    \[
        g \circ f \coloneqq \phi_g \circ (b_g(s_1, \dots, s_n) \circ b_f).
    \]
\end{rem}

We would hope to gain a new functorial definition of a colored operad with the crossed simplicial braid group as its indexing sets. 
\begin{defin}
    A $\Delta B$-category is a category $\Ox$ with a functor
    $$\pi \colon \Ox \to \Delta B$$
    \begin{enumerate}
        \item \textbf{Op-fibration condition:} The functor $\pi$ is an op-fibration.
        \item \textbf{Segal condition:} For $\O \coloneqq \Ox_{\langle 1 \rangle}$, 
        we have $\Ox_{\langle n \rangle} \cong \O^n$ for all $n \geq 1$. 
    \end{enumerate}
\end{defin}

However, this causes some issues. 
We would assume a braided monoidal category to be an example of the new colored operad.
The construction of the braided monoidal structure should be analoguously to the symmetric case.

Meaning, the morphism in $\sigma_0 \in \Delta B$ should induce a braiding defining a braided monoidal structure 
on the underlying category of $\Ox$ using the op-fibration and Segal property.
Thus, similarily to the symmetric case we would want to live $\sigma_0$ using the op-fibration condition.
Hence, the following diagram would have to commute in $\Delta B$:

\begin{center}
    \begin{tikzcd}
        {[1]} \arrow[r, "\tau^1_0"] \arrow[d, "\sigma_0"'] & {[0]} \\
        {[1]} \arrow[ru, "\tau^1_0"']                      &      
    \end{tikzcd}
\end{center}


This is not the case in $\Delta B$. However, we were hopeful that by fixing this we might obtain the 
final candidate for a op-fibration.
So we define a category $\Delta B'$ from $\Delta B$ by requiring for all $n \in \N$ and $1 \leq i \leq n$
\[
    \tau_i^n \circ \sigma_i = \tau_i^n.
\]

There is also a graphical interpretation of $\Delta B$, where $\sigma_i \in B_n$ 
is the usual braid on $n$ strands.
\begin{center}
    \begin{tikzpicture}
            \pic[
            braid/number of strands=6,
            braid/strand 2/.style={white},
            braid/strand 5/.style={white},
            name=coordinates,
            ] at (0,0) {braid={s_3}};
            \node[at=(coordinates-rev-2-e)] {$\dots$};
            \node[at=(coordinates-rev-3-e), below=3pt] {$i$};
            \node[at=(coordinates-rev-4-e), below=3pt] {$i+1$};
            \node[at=(coordinates-rev-5-e)] {$\dots$};
            \node[anchor=center] at (-0.5,-0.75) {$\sigma_i = \quad $};
    \end{tikzpicture}
\end{center}
Order preserving maps in $\Delta$ also have graphical interpretation given by the degeneracy maps
\begin{center}
    \begin{tikzpicture}
                    \draw (-1.5,0)--(-1.5,2);
                    \draw (1.5,0)--(1.5,2);

                    \draw (-0.5,1)--(-0.5,2);
                    \draw (0.5,1)--(0.5,2);
                    \draw (-0.5,1) [bend right=80] to (0.5,1);
                    \draw (0,0.7)--(0,0);



                    \node at (-1.5,2.3) {$1$};
                    \node at (-0.5,2.3) {$i$};
                    \node at (0.5,2.3) {$i+1$};
                    \node at (1.5,2.3) {$n$};
                    \node at (-1.5,-0.3) {$1$};
                    \node at (-0.75, 0) {$\dots$};
                    \node at (0,-0.3) {$i$};
                    \node at (0.75, 0) {$\dots$};
                    \node at (1.5,-0.3) {$n-1$};

                    \node[anchor=center] at (-2.5,1) {$\tau_i^n = \quad $};
    \end{tikzpicture}
\end{center}

and face maps

\begin{center}
        \begin{tikzpicture}
            \pic[
            braid/number of strands=8,
            braid/strand 2/.style={white},
            braid/strand 5/.style={white},
            braid/strand 7/.style={white},
            name=coordinates,
            ] at (0,0) {braid={s_4}};
            %above nodes
            \node[at=(coordinates-rev-1-e), above=1.6] {$1$};
            \node[at=(coordinates-rev-3-e), above=1.6] {$i-1$};
            \node[at=(coordinates-rev-4-e), above=1.6] {$i$};
            \node[at=(coordinates-rev-6-e), above=1.6] {$i+1$};
            \node[at=(coordinates-rev-8-e), above=1.6] {$n$};
            % below nodes
            \node[at=(coordinates-rev-1-e), below=3pt] {$1$};
            \node[at=(coordinates-rev-2-e)] {$\dots$};
            \node[at=(coordinates-rev-3-e), below=3pt] {$i-1$};
            \node[at=(coordinates-rev-4-e), below=3pt] {$i$};
            \node[at=(coordinates-rev-5-e), below=3pt] {$i+1$};
            \node[at=(coordinates-rev-6-e), below=3pt] {$i+2$};
            \node[at=(coordinates-rev-7-e)] {$\dots$};
            \node[at=(coordinates-rev-8-e), below=3pt] {$n+1$};
            \node[anchor=center] at (-0.5,-0.75) {$\delta_i^n = \quad $};
    \end{tikzpicture}
\end{center}

So, we identify the following graphically:

\begin{center}
    \begin{tikzpicture}
            \pic[
            braid/number of strands=6,
            braid/strand 2/.style={white},
            braid/strand 5/.style={white},
            name=coordinates,
            ] at (0,2.5) {braid={s_3}};
            \draw (coordinates-rev-1-e) -- +(0,-1);
            \draw (coordinates-rev-3-e) [bend right=80] to (coordinates-rev-4-e);
            \draw ($ (coordinates-rev-3-e)!0.5!(coordinates-rev-4-e) +(0, -0.3)$) -- +(0,-0.7);
            \draw (coordinates-rev-6-e) -- +(0,-1);

            \pic[
            braid/number of strands=6,
            braid/strand 2/.style={white},
            braid/strand 5/.style={white},
            name=plainbraid,
            ] at (6.5,2.5) {braid={}};
            \draw (plainbraid-rev-1-e) -- +(0,-2);
            \draw (plainbraid-rev-3-e) -- +(0,-1);
            \draw ($ (plainbraid-rev-3-e)+(0, -1)$) [bend right=80] to ($ (plainbraid-rev-4-e) +(0, -1)$);
            \draw ($ (plainbraid-rev-3-e)!0.5!(plainbraid-rev-4-e) +(0, -1.3)$) -- +(0,-0.7);
            \draw (plainbraid-rev-4-e) -- +(0,-1);
            \draw (plainbraid-rev-6-e) -- +(0,-2);

            \node at (5.75, 1) {$=$};


    \end{tikzpicture}
\end{center}
        



As it turns out, a short graphical calculation shows that the hexagon identities hold for $\Delta B'$.x
Unfortunately, we face the next challenge as we discovered that we obtain a braid group action 
where we would not expect one and moreover where we not want a braid group action:

Let $\sigma \in B_n$, then for $\sigma \colon [n] \to [n]$ we obtain cocartesian morphisms in $\Ox_{[n]}$ for
each $X \in \Ox_{[n]} \cong \O^n$, so the braid group $B_n$ acts on $\O^n$. However, we would expect only to
gain a symmetric group action on $\O^n$.

Thus, we have to conclude that the crossed simplicial braid group is not the right candidate for
a functorial definition of a braided operad.